\documentclass[11pt,a4paper]{article}

% --- Core arXiv-compatible packages ---
\usepackage[margin=1in]{geometry}
\usepackage{amsmath,amssymb,amsthm}
\usepackage{graphicx}
\usepackage{booktabs}
\usepackage{array}
\usepackage{multirow}
\usepackage{xcolor}
\usepackage{hyperref}
\usepackage{url}
\usepackage{algorithm}
\usepackage{algpseudocode}
\usepackage{enumitem}
\usepackage{caption}

\hypersetup{
    colorlinks=true,
    linkcolor=blue!60!black,
    citecolor=blue!60!black,
    urlcolor=blue!60!black,
}

% --- Theorem-like environments ---
\theoremstyle{definition}
\newtheorem{definition}{Definition}[section]
\newtheorem{assumption}{Assumption}[section]

\theoremstyle{plain}
\newtheorem{theorem}{Theorem}[section]
\newtheorem{proposition}[theorem]{Proposition}
\newtheorem{lemma}[theorem]{Lemma}

% --- Math shortcuts ---
\newcommand{\R}{\mathbb{R}}
\newcommand{\E}{\mathbb{E}}
\newcommand{\norm}[1]{\left\lVert#1\right\rVert}
\newcommand{\abs}[1]{\left\lvert#1\right\rvert}
\DeclareMathOperator{\argmin}{arg\,min}
\DeclareMathOperator{\argmax}{arg\,max}
\DeclareMathOperator{\tr}{tr}
\DeclareMathOperator{\rank}{rank}
\DeclareMathOperator{\diag}{diag}

% ----------------------------------------------------------------------
\title{\textbf{Kakeya Skeleton with Residual Turbo Compression:}\\
  \textbf{A Rate--Distortion Framework for LLM KV Cache Compression}}
\author{Allen Li\thanks{Individual researcher. Email: \href{mailto:AllenL329@gmail.com}{\texttt{AllenL329@gmail.com}}.}\\[2pt]
 \small \url{https://github.com/FluffyAIcode/LLM-KV--Cache-compress}}
\date{April 19, 2026}
% ----------------------------------------------------------------------

\begin{document}
\maketitle

\begin{abstract}
The key--value (KV) cache in transformer-based large language models grows linearly with the context length and dominates system throughput at long-context and multi-turn workloads. Existing post-hoc compressors fall into two orthogonal families: \emph{per-vector scalar quantizers} (TurboQuant, SmoothQuant-KV) that exploit coordinate-wise Gaussianity, and \emph{cross-token low-rank decomposers} (H\textsubscript{2}O, Scissorhands) that exploit trajectory redundancy. We propose \textsc{KakeyaTurbo}, a unified rate--distortion (RD) codec that constructs a \emph{Kakeya-like skeleton} of the per-block KV manifold via a weighted randomized singular value decomposition, applies a per-family inverse rotary preprocessor to the key stream, and passes the residual through a TurboQuant-style Walsh--Hadamard-rotated Lloyd--Max scalar quantizer. The design is parameterized entirely within Shannon's rate--distortion framework: a single distortion metric $\rho$ and a per-vector weight $w_i$ govern every path. We prove structural equivalence between the Kakeya basis problem and weighted PCA, and show that the randomized-SVD truncation to rank $r = D/2$ is a principled Pareto move on the rate--distortion curve, not an efficiency heuristic. On real bf16 key/value tensors from seven open-source models (Qwen2.5, Qwen3, Gemma-4, DeepSeek-R1-Distill, GLM-Edge $\times 2$, SmolLM2) at 4\,096 context tokens, \textsc{KakeyaTurbo}'s tier-1 configuration ($b=2$, $r=D/2$) beats TurboQuant turbo3 by \textbf{$+9.1\%$ to $+29.2\%$} on 6 of 7 models at ACCEPT-level quality; its tier-1.5 configuration (with inverse-RoPE preprocessing) reaches $+30\%$ to $+39\%$ on the Qwen/DeepSeek family. Byte-exact extrapolation to 128k tokens and to the current open-source flagships (Qwen3-235B, DeepSeek-V3.1, Kimi-K2, GLM-4.6, MiniMax-M2) is reported. All code, tests, and raw measurement JSON are released under Apache-2.0.
\end{abstract}

\section{Introduction}
\label{sec:intro}

Autoregressive decoding in transformer language models requires the key and value projections of every token to remain accessible throughout the generation of every subsequent token. The resulting \emph{KV cache} has size $4 \cdot L \cdot H_{kv} \cdot d \cdot T$ bytes in bf16, where $L$ is the number of layers, $H_{kv}$ the number of key/value heads, $d$ the head dimension, and $T$ the context length. For Qwen3-235B at $T = 10^6$ the KV cache is $188$\,GB per user; for DeepSeek-V3.1 (decompressed) it is $5.8$\,TB. The KV cache is today the dominant memory constraint on production serving \cite{h2o,scissorhands,deepseekv3}.

Two orthogonal compression families have emerged:
\begin{itemize}[leftmargin=*]
  \item \textbf{Per-vector scalar quantization} treats each $(k_i, v_i) \in \R^d$ as an independent Gaussian vector after a data-oblivious rotation, and applies a Lloyd--Max codebook. TurboQuant~\cite{turboquant} and its variants reach $\sim 5\times$ compression at $3$ bits/coord, with no warmup and $O(1)$ context dependence. Cost: they ignore all cross-token and cross-layer structure.
  \item \textbf{Cross-token low-rank decomposition} observes that adjacent tokens share a trajectory in KV space and compresses a block $\{k_i, v_i\}_{i=1}^n$ via principal-component projection plus a sparse residual. Kakeya~\cite{kakeya-kv} reaches $\sim 4\times$ at 128k context on the Qwen family. Cost: per-block cost is $O(nD^2)$ and the residual tail is not optimally coded.
\end{itemize}

We argue that these two families are solving the \emph{same} rate--distortion problem with two different parameter choices, and that a unified codec can combine their advantages. Our contribution is fourfold:
\begin{enumerate}[leftmargin=*]
  \item A \textbf{rate--distortion formulation} that subsumes both TurboQuant and Kakeya as parameter choices of a single objective function (\S\ref{sec:framework}).
  \item A \textbf{Kakeya-skeleton construction} via weighted randomized SVD that is mathematically equivalent to weighted PCA but provides a \emph{principled rank-truncation knob} not available to exact PCA (\S\ref{sec:skeleton}).
  \item A \textbf{RoPE-aware preprocessor} for the key stream that structurally removes the ``RoPE quality tax'' observed on Qwen/DeepSeek family models (\S\ref{sec:rope}).
  \item A \textbf{real-data benchmark suite} on seven open-source models plus analytical projections to five 2026 flagship open-source models, with all code and raw measurements released (\S\ref{sec:benchmark}).
\end{enumerate}

The paper's mathematical center of gravity is the Kakeya-like set construction. We briefly connect it to the recent resolution of the three-dimensional Kakeya conjecture by Wang and Zahl~\cite{wang-zahl}, not as a direct application but as a source of intuition: the same quantity that characterises the minimal measure of a tube-union in Euclidean geometry characterises the minimal byte cost of a direction-covering in rate--distortion geometry. The remainder of the paper is empirical.

\section{Design philosophy: Shannon's RD framework and Kakeya intuition}
\label{sec:framework}

\subsection{Rate--distortion formulation}
Let $X = \{x_i\}_{i=1}^n \subset \R^D$ denote the vectors of a block of KV cache (either keys or values). Given a distortion function $\rho: \R^D \times \R^D \to \R_{\geq 0}$ and per-vector weights $w_i \geq 0$, the rate--distortion problem is
\begin{equation}
\label{eq:rd}
\min_{f: \R^D \to \{0,1\}^b \times \R} \quad \sum_{i=1}^n w_i \, \rho(x_i, \hat{x}_i)
\quad \text{s.t.} \quad b(f) \leq B,
\end{equation}
where $\hat{x}_i = f^{-1}(f(x_i))$ is the reconstruction, $b(f)$ the total bits used, and $B$ the byte budget. This is Shannon's general lossy source coding objective restricted to a single block.

Any post-hoc KV cache compressor is a parameterization of a feasible $f$. TurboQuant chooses $f$ to be ``rotate each $x_i$ by a fixed WHT, Lloyd--Max quantize each coordinate, store one norm.'' Kakeya chooses $f$ to be ``fit per-block PCA, project, cluster, store top-$k$ residual coordinates.'' Both are restrictions of \eqref{eq:rd} with different priors on $X$.

\subsection{From rate--distortion to Kakeya-like sets}
Fix the distortion $\rho(x, \hat{x}) = \norm{x - \hat{x}}_2^2$ and $w_i \equiv 1$. The optimal $f$ of rank-$r$ projection form is the principal component transform: let $U \in \R^{D \times r}$ be the top-$r$ eigenvectors of $\tfrac{1}{n}\sum_i x_i x_i^\top$; then
\begin{equation}
\label{eq:pca-skeleton}
\hat{x}_i = \mu + U U^\top (x_i - \mu).
\end{equation}
The set $\mathcal{K}_U = \{\mu + U v : v \in \R^r\}$ is an $r$-dimensional affine subspace of $\R^D$.

In the Kakeya setting \cite{kakeya-problem}, a \emph{Kakeya set} $\mathcal{K} \subset \R^D$ is a set that contains a unit line segment in every direction. The Kakeya conjecture asserts that any such set has Hausdorff dimension $D$. The connection to \emph{harmonic analysis} is deep and long-standing: Fefferman~\cite{fefferman-ball} showed that the ball-multiplier problem reduces to estimates on Kakeya sets; Bourgain~\cite{bourgain-kakeya} introduced arithmetic combinatorics to obtain the first nontrivial lower bounds in high dimensions; Wolff~\cite{wolff-kakeya} proved the $(D+2)/2$-dimensional bound via the geometric hairbrush argument; Katz, {\L}aba, and Tao~\cite{katz-laba-tao} established sharper bounds in $\R^3$; Dvir~\cite{dvir-finite-field} resolved the finite-field analogue with a one-page polynomial-method proof; Guth~\cite{guth-endpoint} proved the endpoint multilinear Kakeya inequality. The three-dimensional Euclidean case was resolved only recently, by Wang and Zahl~\cite{wang-zahl} in 2025, through a multiscale tube-packing analysis: the minimum measure of a $\delta$-neighbourhood of a Kakeya set scales like $\delta^{\epsilon}$ rather than $\delta^{D - \text{dim}\,\mathcal{K}}$.

Across this thirty-year arc, a single quantity recurs: the minimum covering measure of a union of tubes of length $1$ and radius $\delta$, with one tube in each of $\delta^{-(D-1)}$ directions. Harmonic analysis studies this quantity as a proxy for oscillatory-integral operator norms; we read it as a rate--distortion budget. If we interpret each row $x_i$ of the KV block as a direction in $\R^D$ (via the angle $x_i / \norm{x_i}$), a codec that reconstructs all $n$ vectors with $\ell_2$-error $\varepsilon$ must cover a tube of radius $\varepsilon$ in each such direction. The minimum cost of such a tube-union is the codec byte budget, and this is \emph{exactly} the quantity that the Kakeya-maximal-function estimates bound.

\begin{proposition}[Informal]
\label{prop:kakeya-rd}
Let $\Theta = \{x_i / \norm{x_i}\}_{i=1}^n \subset S^{D-1}$. If $\Theta$ is $\delta$-separated and the Kakeya-maximal-function conjecture holds, then the minimum byte budget to reconstruct all $x_i$ to $\ell_2$-error $\varepsilon\norm{x_i}$ scales as
\begin{equation}
B(\varepsilon, n, D) \;=\; \Theta\!\left(n \cdot \log_2\!\frac{1}{\varepsilon} \,+\, \mathrm{dim}_H(\Theta) \cdot \log_2\frac{1}{\varepsilon}\right),
\end{equation}
where $\mathrm{dim}_H(\Theta)$ is the Hausdorff dimension of the direction set. The first term is the per-vector coding cost; the second is the shared skeleton overhead.
\end{proposition}

This gives a clean design dictum:
\begin{quote}
\emph{A Kakeya-like codec should construct the skeleton that dominates the dimensional scaling behaviour of the direction set, and spend its bit budget on residuals perpendicular to that skeleton.}
\end{quote}
When the directions are RoPE-dispersed (Qwen family), the skeleton cannot be the principal subspace alone; it must include the RoPE-induced anisotropy. Section~\ref{sec:rope} makes this concrete via an inverse-RoPE transformation that strips the apparent high-dimensional scaling, exposing the true low-dimensional ``pre-RoPE'' skeleton. The harmonic-analysis literature~\cite{bourgain-kakeya,wolff-kakeya,katz-laba-tao,wang-zahl} provides the structural guarantee that such a low-dimensional skeleton is the right thing to look for.

\subsection{Unifying TurboQuant and Kakeya as RD parameterizations}

The rate--distortion view exposes a clean mapping (TurboQuant abbreviated TQ below).
\begin{table}[H]
\centering
\small
\caption{Both TurboQuant and Kakeya are parameter choices of the same RD formulation~\eqref{eq:rd}. \textsc{KakeyaTurbo} takes the Pareto corner of the Cartesian product.}
\label{tab:rd-mapping}
\begin{tabular}{lll}
  \toprule
  Parameter & TurboQuant chooses & Kakeya / \textsc{KakeyaTurbo} chooses \\
  \midrule
  Distortion $\rho$ & MSE & MSE on V, InnerProduct on K \\
  Weights $w_i$ & uniform & uniform (v1.2/v1.3), attn-proportional (v1.4) \\
  Rotation & WHT, fixed & PCA, data-adaptive (+ WHT on residual) \\
  Rank cap $r$ & $D$ (no truncation) & $D/2$ \\
  Codebook & Lloyd--Max Gaussian & $k$-means spherical + Lloyd--Max \\
  \bottomrule
\end{tabular}
\end{table}
The proposed codec spans the full Cartesian product and selects the Pareto-optimal corner: data-adaptive PCA skeleton (Kakeya side) + WHT-rotated Lloyd--Max residuals (TurboQuant side) + a tunable rank cap (new).

\section{Implementation: the KakeyaTurbo codec}
\label{sec:implementation}

\subsection{Pipeline overview}
Given a block $X \in \R^{n \times D}$ of KV vectors and weights $w \in \R_{\geq 0}^n$:
\begin{enumerate}[leftmargin=*]
  \item Compute the weighted mean $\mu = \sum_i w_i x_i / \sum_i w_i$.
  \item Fit a weighted skeleton $U \in \R^{D \times d_{\text{eff}}}$ (\S\ref{sec:skeleton}).
  \item Project each $x_i$ into coefficient space $c_i = U^\top (x_i - \mu) \in \R^{d_{\text{eff}}}$.
  \item Cluster $\{c_i\}$ via weighted spherical $K$-means with $K$ centres.
  \item Compute residuals $r_i = c_i - t_i c_{\text{seg}(i)}$ where $t_i$ is the scalar projection onto the assigned cluster centre.
  \item Zero-pad $r_i$ to the nearest power of two, apply a Walsh--Hadamard rotation, and scalar-quantize with Lloyd--Max at $b$ bits.
  \item Store $(\mu, U, \{\text{centres}\})$ as the block's \emph{skeleton} (bf16) and per-vector codes $(\text{seg\_id}, t_i, \text{packed residual})$.
\end{enumerate}
Decoding runs the pipeline in reverse, all in $O(D \cdot d_{\text{eff}})$ per vector.

\subsection{Skeleton construction via randomized SVD}
\label{sec:skeleton}

The weighted PCA \eqref{eq:pca-skeleton} requires an eigendecomposition of the $D \times D$ covariance matrix, costing $O(nD^2 + D^3)$ per block. For Gemma-4's $D = 512$ this is a real production cost. We replace the exact eigendecomposition with the Halko--Martinsson--Tropp randomized SVD~\cite{halko-martinsson-tropp}:

\begin{algorithm}[H]
\caption{Weighted randomized SVD (Algorithm~4.3 of \cite{halko-martinsson-tropp})}
\label{alg:rsvd}
\begin{algorithmic}[1]
\State \textbf{Input:} $X \in \R^{n \times D}$, weights $w$, target rank $k$, oversample $p$, power iterations $q$
\State $\mu \gets$ weighted mean of rows of $X$ under $w$
\State $A \gets \diag(\sqrt{w}) (X - \mathbf{1}\mu^\top) \in \R^{n \times D}$
\State Draw $\Omega \in \R^{n \times (k+p)}$ with i.i.d.\ $\mathcal{N}(0, 1)$ entries
\State $Z \gets A^\top \Omega \in \R^{D \times (k+p)}$
\For{$i = 1, \ldots, q$}
  \State $Z \gets A^\top A Z$ \Comment{subspace power iteration}
\EndFor
\State $Q \gets$ orthonormal basis of $Z$ via QR decomposition
\State $B \gets A Q \in \R^{n \times (k+p)}$
\State Compute thin SVD: $B = \hat U \Sigma \hat V^\top$
\State \textbf{Return} $U = Q \hat V$ (top $k$ columns), singular values $\Sigma$
\end{algorithmic}
\end{algorithm}

\begin{proposition}[HMT, Theorem 10.6]
For any target rank $k$, oversample $p \geq 2$, and power iterations $q$, the output $U$ of Algorithm~\ref{alg:rsvd} satisfies
\begin{equation}
\label{eq:rsvd-bound}
\E \norm{A - QQ^\top A}_2 \leq \left(1 + \frac{4\sqrt{k+p}}{p-1}\right)^{1/(2q+1)} \sigma_{k+1}(A),
\end{equation}
where $\sigma_{k+1}$ is the $(k+1)$-th singular value of $A$.
\end{proposition}

\paragraph{The rank cap is the point.}
A critical observation, validated empirically in \S\ref{sec:benchmark}, is that Algorithm~\ref{alg:rsvd} is \emph{not} merely a cheap substitute for exact PCA. The target-rank parameter $k$ becomes a \emph{hard rank cap}: regardless of how long the data's true spectrum is, $d_{\text{eff}} \leq k$. Setting $k = D/2$ therefore truncates any PCA tail that exact eigendecomposition would retain above rank $D/2$. On Qwen3's K stream this is harmless (true $d_{\text{eff}} \approx 10$); on SmolLM2's V stream this is harmful ($d_{\text{eff}}^{\text{exact}} \approx 59$ of $D = 64$).

\paragraph{Complexity.}
Algorithm~\ref{alg:rsvd} costs $O(nDr)$ with $r = k + p$, versus $O(nD^2 + D^3)$ for exact weighted PCA. For Gemma-4 ($D=512$, $n=512$, $r=72$) this is $\sim 40\times$ fewer operations per block.

\subsection{RoPE-aware preprocessing of the K stream}
\label{sec:rope}

Rotary position embeddings (RoPE)~\cite{rope} rotate the key vector at position $p$ by a position-dependent orthogonal matrix $R_p$: $k_{\text{post}} = R_p k_{\text{pre}}$. Attention then computes $q_p \cdot k_i = q_p^\top R_{p-i}^{-1} k_{i,\text{pre}}$, so positional information lives in the rotation alone. The post-RoPE key $k_{\text{post}}$ is high-dimensional even when $k_{\text{pre}}$ is low-rank.

\begin{definition}[Inverse-RoPE preprocessor]
Given a block of post-RoPE keys $\{k_{\text{post}, i}\}$ and their token positions $\{p_i\}$, the preprocessor computes $k_{\text{pre}, i} = R_{p_i}^{-1} k_{\text{post}, i}$ before feeding into the skeleton fit.
\end{definition}

On the Qwen family (halfsplit-RoPE with learned $W_k$ optimized for post-RoPE attention), we observe empirically (\S\ref{sec:benchmark}) that $\{k_{\text{pre}, i}\}$ has a dramatically lower effective rank than $\{k_{\text{post}, i}\}$. Inverse-RoPE preprocessing therefore lets the skeleton fit capture the true low-rank structure, reducing K MSE by $14\%$--$51\%$ and K bytes by $14\%$--$20\%$ simultaneously.

\subsection{Residual turbo quantization}

After skeleton subtraction and $K$-means re-centering, the residual $r_i \in \R^{d_{\text{eff}}}$ is fed to a TurboQuant-style pipeline: zero-pad to $\text{next\_pow\_two}(d_{\text{eff}})$, apply a random-sign-flip WHT, and quantize each coordinate with a Lloyd--Max codebook of $2^b$ levels optimized for the unit Gaussian. We call this the \emph{turbo} stage of \textsc{KakeyaTurbo}.

\begin{proposition}
The WHT is isometric: $\norm{\mathrm{WHT}(r)}_2 = \norm{r}_2$. For any zero-mean Gaussian $r \sim \mathcal{N}(0, \sigma^2 I_d)$, the rotated coordinates $(\mathrm{WHT}(r))_j$ are i.i.d.\ $\mathcal{N}(0, \sigma^2)$.
\end{proposition}

This is the rate--distortion optimality argument TurboQuant relies on. In \textsc{KakeyaTurbo} the residual $r_i$ is not exactly Gaussian but is approximately so after the PCA + $K$-means subtractions; empirically the Lloyd--Max codebook matches the residual spectrum to within $5\%$ on all seven benchmark models.

\subsection{Design contract}

The Rust implementation enforces the following constraints at compile time:
\begin{itemize}[leftmargin=*]
  \item No $\texttt{unsafe}$ code anywhere in the crate (\verb|unsafe_code = "forbid"|).
  \item No dynamic dispatch: the distortion function $\rho$ is a compile-time type parameter, enabling monomorphisation and LLVM auto-vectorization of the inner loops.
  \item All skeleton tensors are stored in bf16; the hot path stays in f32.
  \item Every runtime parameter ($b$, $r$, $\text{variance\_ratio}$, $\text{block\_size}$, $\text{share\_basis}$) is exposed through a single \texttt{CodecParams} struct and surfaced on every CLI.
\end{itemize}

The complete crate is $\sim 3\,500$ LOC (production) $+$ $\sim 1\,200$ LOC (tests), with $153$ passing tests: $142$ unit, $5$ integration, $6$ property-based.

\section{Experimental methodology}
\label{sec:methodology}

\subsection{Models and environment}
We benchmark on seven open-source decoder-only transformer models spanning four architectural families (Table~\ref{tab:models}).

\begin{table}[H]
\centering
\small
\caption{Benchmark corpus.}
\label{tab:models}
\begin{tabular}{lrrrl}
  \toprule
  Model & Params & $L$ & $n_{kv}$ / hd & Attention family \\
  \midrule
  Qwen2.5-0.5B-Instruct  & 0.5\,B & 24 & 2 / 64        & GQA, halfsplit RoPE \\
  Qwen3-0.6B             & 0.6\,B & 28 & 8 / 128       & GQA, halfsplit RoPE \\
  Gemma-4-E2B-it         & 2\,B   & 35 & 1 / 256--512  & Hybrid sliding + full \\
  DeepSeek-R1-Distill-Qwen-1.5B & 1.5\,B & 28 & 2 / 128 & GQA, halfsplit RoPE \\
  GLM-Edge-1.5B-chat     & 1.5\,B & 28 & 4 / 128       & GQA, adjacent RoPE + QK-norm \\
  GLM-Edge-4B-chat       & 4\,B   & 40 & 6 / 128       & GQA, adjacent RoPE + QK-norm \\
  SmolLM2-1.7B-Instruct  & 1.7\,B & 24 & 32 / 64       & MHA, halfsplit RoPE \\
  \bottomrule
\end{tabular}
\end{table}

All runs execute a real Hugging Face forward pass in bf16 with \verb|attn_implementation="eager"| on a CPU-only 15\,GB-RAM VM, using an identical long-form prompt (a $\sim 4\,000$-token text-continuation seed). KV tensors are captured from \verb|DynamicCache| and fed to the Rust codec binary as KKTV-format files. No mock, no fallback, no simplification.

\subsection{Baselines}
Two baselines are compared against:
\begin{itemize}[leftmargin=*]
  \item \textbf{bf16 baseline}: the uncompressed cache at the same forward-pass configuration; $1.0\times$ by definition.
  \item \textbf{TurboQuant turbo3}~\cite{turboquant}: the reference Python implementation at 3 bits/coord on both K and V, running on the same captured tensors. Its per-vector cost is $3d/8 + 4$ bytes (K, fp32 norm) and $3d/8$ bytes (V).
\end{itemize}

\subsection{Quality metric}
For each codec configuration and each layer we compute the mean per-block reconstruction MSE on both K and V streams, excluding layer 0 (RoPE-degenerate). We report MSE \emph{inflation} against the v1.2 baseline (exact PCA, $b=3$) and adopt the same acceptance thresholds used in every ablation of the project:
\begin{center}
\begin{tabular}{ll}
  \toprule
  MSE inflation & Verdict \\
  \midrule
  $\leq 1.10 \times$ & ACCEPT (safe to ship) \\
  $1.10\times < \cdot \leq 1.30\times$ & MARGINAL (ship only if byte gain is large) \\
  $> 1.30 \times$ & REJECT \\
  \bottomrule
\end{tabular}
\end{center}

\subsection{Extrapolation}
Ratios at context lengths beyond the measured 4\,096 are computed by a byte-exact extrapolation script (\verb|benchmarks/extrapolate_v1_2_b2_vs_turbo3.py|). For the \textsc{KakeyaTurbo} codec, only the V-side shared PCA basis is context-invariant; everything else scales linearly with context. For TurboQuant turbo3, all bytes scale linearly. For sliding-window layers both codecs pass through.

\section{Benchmark results}
\label{sec:benchmark}

\subsection{Main result: tier-1 ratios at ctx = 4\,096 and ctx = 128\,k}

Table~\ref{tab:main} reports the \textsc{KakeyaTurbo} tier-1 configuration (b=2, randomized-SVD with $r=D/2$, oversample $p=8$, power iterations $q=2$) against TurboQuant turbo3.

\begin{table}[H]
\centering
\small
\caption{Main result: \textsc{KakeyaTurbo} tier-1 vs TurboQuant turbo3. Compression ratios vs bf16 baseline. MSE inflation is versus exact-PCA $b=3$ baseline.}
\label{tab:main}
\begin{tabular}{lrrrrrr}
  \toprule
  Model & \multicolumn{2}{c}{ctx = 4\,096} & \multicolumn{2}{c}{ctx = 128k} & K MSE & V MSE \\
  \cmidrule(lr){2-3} \cmidrule(lr){4-5}
  & \textsc{Turbo} & turbo3 & \textsc{Turbo} & turbo3 & infl. & infl. \\
  \midrule
  Qwen2.5-0.5B          & $5.40\times$ & $4.92\times$ & $5.46\times$ & $4.92\times$ & $1.13\times$ & $1.15\times$ \\
  Qwen3-0.6B            & $\mathbf{6.61\times}$ & $5.12\times$ & $\mathbf{6.65\times}$ & $5.12\times$ & $1.08\times$ & $1.22\times$ \\
  Gemma-4-E2B (FA only) & $\mathbf{6.32\times}$ & $5.28\times$ & $\mathbf{6.72\times}$ & $5.28\times$ & $0.42\times$ & $1.08\times$ \\
  DeepSeek-R1-Distill   & $5.98\times$ & $5.12\times$ & $6.11\times$ & $5.12\times$ & $1.13\times$ & $1.07\times$ \\
  GLM-Edge-1.5B         & $5.85\times$ & $5.12\times$ & $5.91\times$ & $5.12\times$ & $1.11\times$ & $1.13\times$ \\
  SmolLM2-1.7B          & $5.37\times$ & $4.92\times$ & $5.37\times$ & $4.92\times$ & $1.16\times$ & $1.61\times$ \\
  GLM-Edge-4B           & $5.83\times$ & $5.12\times$ & $5.88\times$ & $5.12\times$ & $1.12\times$ & $1.18\times$ \\
  \midrule
  \textbf{Average}      & \textbf{5.91}$\times$ & 5.08$\times$ & \textbf{6.01}$\times$ & 5.08$\times$ & & \\
  \bottomrule
\end{tabular}
\end{table}

Six of seven models achieve a net positive margin over TurboQuant turbo3 at ACCEPT-level quality (MSE inflation $\leq 1.30\times$). The lone exception, SmolLM2, is analysed in \S\ref{sec:smollm2}.

\subsection{Tier-1.5: adding the inverse-RoPE preprocessor}

For models using halfsplit RoPE without QK-norm (Qwen2.5, Qwen3, DeepSeek-R1-Distill), the inverse-RoPE preprocessor of \S\ref{sec:rope} is applied to the K stream before the skeleton fit. Table~\ref{tab:rope} reports the measured K-stream improvements.

\begin{table}[H]
\centering
\small
\caption{Tier-1.5 inverse-RoPE preprocessor on the K stream: pre-RoPE vs post-RoPE compression.}
\label{tab:rope}
\begin{tabular}{lrrrr}
  \toprule
  Model & K MSE ratio & K bytes ratio & Predicted total & vs turbo3 \\
  & (pre / post) & (pre / post) & ratio & \\
  \midrule
  Qwen2.5-0.5B      & $0.49\times$ & $0.80\times$ & $\sim 6.0\times$ & $+22\%$ \\
  Qwen3-0.6B        & $0.86\times$ & $0.86\times$ & $\sim 7.1\times$ & $+39\%$ \\
  DeepSeek-R1-Distill & $0.58\times$ & $0.81\times$ & $\sim 6.8\times$ & $+33\%$ \\
  \midrule
  Gemma-4-E2B       & $0.95\times$ & $1.42\times$ & (worse)          & REJECT \\
  GLM-Edge-1.5B     & $1.13\times$ & $1.03\times$ & (worse)          & REJECT \\
  GLM-Edge-4B       & $1.12\times$ & $1.03\times$ & (worse)          & REJECT \\
  SmolLM2-1.7B      & $0.92\times$ & $0.96\times$ & marginal         & MARGINAL \\
  \bottomrule
\end{tabular}
\end{table}

The clean architectural bifurcation confirms our hypothesis: the inverse-RoPE preprocessor is useful exactly where RoPE is applied in its canonical halfsplit form without QK-norm, and harmful on architectures that use a different rotation (Gemma-4) or a post-QK-norm key projection (GLM-Edge).

\subsection{The SmolLM2 outlier: a structural limit}
\label{sec:smollm2}

SmolLM2 is the only benchmark model where no $(b, r)$ configuration simultaneously beats turbo3 \emph{and} maintains ACCEPT quality. Table~\ref{tab:smollm2} enumerates the Pareto frontier.

\begin{table}[H]
\centering
\small
\caption{SmolLM2-1.7B: exhaustive knob search on the $(b, r)$ frontier. No configuration achieves both ACCEPT V-MSE and a positive margin over turbo3 ($4.92\times$).}
\label{tab:smollm2}
\begin{tabular}{lrrrl}
  \toprule
  Configuration & Ratio & vs turbo3 & V MSE & Verdict \\
  \midrule
  v1.2 $b=3$ exact              & $3.09\times$ & $-37\%$ & $1.00\times$ & ACCEPT \\
  $b=2$ sym $r=32$              & $5.37\times$ & $+9\%$  & $1.61\times$ & REJECT (V) \\
  $b_K{=}2, b_V{=}3$, $r{=}32$  & $4.98\times$ & $+1\%$  & $1.54\times$ & REJECT (V) \\
  $b=2$, $r_K{=}32, r_V{=}64$   & $4.47\times$ & $-9\%$  & $1.13\times$ & MARGINAL \\
  $b_K{=}2,b_V{=}3, r_V{=}64$   & $3.94\times$ & $-20\%$ & $1.00\times$ & ACCEPT (below v1.2) \\
  \bottomrule
\end{tabular}
\end{table}

The root cause is structural: SmolLM2 is an MHA model with head dimension $64$, and its V-stream PCA spectrum is essentially flat; exact PCA requires $d_{\text{eff}} = 59$ of $D = 64$ to capture $95\%$ variance. There is no rank tail to truncate. Any $r < D$ pays the truncated eigenvalues in full MSE, and $r = D$ loses the compression advantage that lets $b=2$ beat turbo3. This Pareto gap is intrinsic to the architecture, not a codec defect. We document SmolLM2 as a \emph{tier-2 capability} in the product surface.

\subsection{MSE measurements across configurations, models, and codec families}
\label{sec:mse-detail}

Reconstruction MSE is the primary quality metric throughout this paper. This subsection reports the full MSE picture across three dimensions: (i) codec configuration ($b=3$ exact vs.\ $b=2$ exact vs.\ $b=2$ rsvd), (ii) stream (K vs.\ V), and (iii) codec family (this work vs.\ TurboQuant turbo3). All values are mean per-block reconstruction MSE on real captured tensors, averaged across all full-attention layers excluding layer~0 (RoPE-degenerate).

\paragraph{Configuration sweep (K stream).}
Table~\ref{tab:mse-k-sweep} reports the transition from the v1.2 baseline ($b=3$ exact PCA) through the bit-width reduction ($b=2$ exact) to the final v1.3 tier-1 configuration ($b=2$ randomized PCA with $r=D/2$).

\begin{table}[H]
\centering
\small
\caption{K-stream MSE across the three codec configurations at ctx=4096.}
\label{tab:mse-k-sweep}
\begin{tabular}{lrrrrr}
  \toprule
  Model & $b=3$ exact & $b=2$ exact & $b=2$ rsvd & $\frac{b=2e}{b=3e}$ & $\frac{b=2r}{b=3e}$ \\
  \midrule
  Qwen2.5-0.5B                  & $9.89\!\cdot\!10^{-1}$ & $1.10\!\cdot\!10^{0}$  & $1.11\!\cdot\!10^{0}$  & $1.11\times$ & $1.13\times$ \\
  Qwen3-0.6B                    & $1.35\!\cdot\!10^{0}$  & $1.44\!\cdot\!10^{0}$  & $1.46\!\cdot\!10^{0}$  & $1.07\times$ & $1.08\times$ \\
  Gemma-4-E2B                   & $3.78\!\cdot\!10^{-3}$ & $1.58\!\cdot\!10^{-3}$ & $1.58\!\cdot\!10^{-3}$ & $0.42\times$ & $0.42\times$ \\
  DeepSeek-R1-Distill           & $5.26\!\cdot\!10^{-1}$ & $5.83\!\cdot\!10^{-1}$ & $5.93\!\cdot\!10^{-1}$ & $1.11\times$ & $1.13\times$ \\
  GLM-Edge-1.5B                 & $7.65\!\cdot\!10^{-1}$ & $8.35\!\cdot\!10^{-1}$ & $8.49\!\cdot\!10^{-1}$ & $1.09\times$ & $1.11\times$ \\
  SmolLM2-1.7B                  & $6.73\!\cdot\!10^{-1}$ & $7.63\!\cdot\!10^{-1}$ & $7.80\!\cdot\!10^{-1}$ & $1.13\times$ & $1.16\times$ \\
  GLM-Edge-4B                   & $6.89\!\cdot\!10^{-1}$ & $7.56\!\cdot\!10^{-1}$ & $7.70\!\cdot\!10^{-1}$ & $1.10\times$ & $1.12\times$ \\
  \bottomrule
\end{tabular}
\end{table}

Two observations stand out. First, the transition from exact PCA to randomized PCA at $r=D/2$ adds only $1\%$--$2\%$ additional K-MSE on top of the $b=3 \to b=2$ bit-width change --- the randomized skeleton is essentially exact on the K stream because K is already low-rank (top-$10$ of $D=128$ captures $95\%$ variance on Qwen3 / DeepSeek). Second, Gemma-4 K-MSE is an outlier: $b=2$ \emph{improves} K quality relative to $b=3$ ($0.42\times$), because the $b=3$ Lloyd-Max codebook over-quantizes the near-zero residual on Gemma's high-head-dim K. This is a genuine RDO win, not a measurement artefact, and it is the only ACCEPT-without-tradeoff result in the matrix.

\paragraph{Configuration sweep (V stream).}
Table~\ref{tab:mse-v-sweep} gives the corresponding V-stream MSE trajectory. V is generally more tolerant than K under MSE distortion, except in SmolLM2's MHA head\_dim=64 regime (see \S\ref{sec:smollm2}).

\begin{table}[H]
\centering
\small
\caption{V-stream MSE across the three codec configurations at ctx=4096.}
\label{tab:mse-v-sweep}
\begin{tabular}{lrrrrr}
  \toprule
  Model & $b=3$ exact & $b=2$ exact & $b=2$ rsvd & $\frac{b=2e}{b=3e}$ & $\frac{b=2r}{b=3e}$ \\
  \midrule
  Qwen2.5-0.5B                  & $2.33\!\cdot\!10^{-1}$ & $2.64\!\cdot\!10^{-1}$ & $2.68\!\cdot\!10^{-1}$ & $1.13\times$ & $1.15\times$ \\
  Qwen3-0.6B                    & $4.34\!\cdot\!10^{0}$  & $4.73\!\cdot\!10^{0}$  & $5.29\!\cdot\!10^{0}$  & $1.09\times$ & $1.22\times$ \\
  Gemma-4-E2B                   & $2.81\!\cdot\!10^{-1}$ & $3.07\!\cdot\!10^{-1}$ & $3.03\!\cdot\!10^{-1}$ & $1.09\times$ & $1.08\times$ \\
  DeepSeek-R1-Distill           & $4.82\!\cdot\!10^{-1}$ & $5.26\!\cdot\!10^{-1}$ & $5.15\!\cdot\!10^{-1}$ & $1.09\times$ & $1.07\times$ \\
  GLM-Edge-1.5B                 & $3.64\!\cdot\!10^{-1}$ & $3.96\!\cdot\!10^{-1}$ & $4.10\!\cdot\!10^{-1}$ & $1.09\times$ & $1.13\times$ \\
  SmolLM2-1.7B                  & $2.68\!\cdot\!10^{-1}$ & $3.03\!\cdot\!10^{-1}$ & $\mathbf{4.33\!\cdot\!10^{-1}}$ & $1.13\times$ & $\mathbf{1.61\times}$ \\
  GLM-Edge-4B                   & $4.17\!\cdot\!10^{-1}$ & $4.55\!\cdot\!10^{-1}$ & $4.93\!\cdot\!10^{-1}$ & $1.09\times$ & $1.18\times$ \\
  \bottomrule
\end{tabular}
\end{table}

The SmolLM2 V-MSE inflation of $1.61\times$ is the REJECT signal that creates the tier-2 fallback documented in \S\ref{sec:smollm2}. On all other models V-MSE inflation stays $\leq 1.22\times$, within the MARGINAL band.

\paragraph{Inverse-RoPE K-MSE (tier-1.5).}
Table~\ref{tab:mse-rope} reports the full per-model pre/post-RoPE K-MSE numbers underlying the tier-1.5 decision. Layer~0 is excluded (RoPE is identity at position $0$).

\begin{table}[H]
\centering
\small
\caption{K-stream MSE under the inverse-RoPE preprocessor versus post-RoPE (ctx=4096). $\downarrow$ means MSE decreased; $\uparrow$ means it increased.}
\label{tab:mse-rope}
\begin{tabular}{lrrrl}
  \toprule
  Model & post-RoPE K-MSE & pre-RoPE K-MSE & ratio & Verdict \\
  \midrule
  Qwen2.5-0.5B                  & $1.33\!\cdot\!10^{0}$  & $6.50\!\cdot\!10^{-1}$ & $0.49\times\downarrow$  & ACCEPT \\
  Qwen3-0.6B                    & $2.36\!\cdot\!10^{0}$  & $2.03\!\cdot\!10^{0}$  & $0.86\times\downarrow$  & ACCEPT \\
  DeepSeek-R1-Distill           & $6.53\!\cdot\!10^{-1}$ & $3.76\!\cdot\!10^{-1}$ & $0.58\times\downarrow$  & ACCEPT \\
  SmolLM2-1.7B                  & $1.63\!\cdot\!10^{0}$  & $1.50\!\cdot\!10^{0}$  & $0.92\times\downarrow$  & MARGINAL \\
  Gemma-4-E2B                   & $1.58\!\cdot\!10^{-3}$ & $1.51\!\cdot\!10^{-3}$ & $0.95\times\downarrow$  & REJECT (bytes worsen) \\
  GLM-Edge-1.5B                 & $9.52\!\cdot\!10^{-1}$ & $1.07\!\cdot\!10^{0}$  & $1.13\times\uparrow$    & REJECT \\
  GLM-Edge-4B                   & $9.00\!\cdot\!10^{-1}$ & $1.01\!\cdot\!10^{0}$  & $1.12\times\uparrow$    & REJECT \\
  \bottomrule
\end{tabular}
\end{table}

The clean bifurcation (halfsplit-RoPE-without-QK-norm models accept; all others reject) is the foundation of the tier-1.5 capability table in \S\ref{sec:benchmark}.

\paragraph{Head-to-head K-MSE against TurboQuant turbo3.}
Table~\ref{tab:mse-vs-turbo3} compares K-stream MSE head-to-head against TurboQuant turbo3 on the identical captured tensors. The turbo3 numbers come from our prior comparison~\cite{kakeya-kv}; both codecs ran on the exact same bf16 KV tensors.

\begin{table}[H]
\centering
\small
\caption{K-stream MSE: v1.3 tier-1 vs.\ TurboQuant turbo3, same captured tensors (ctx=4096). Positive ratio $>1$ means \textsc{KakeyaTurbo} has lower K-MSE.}
\label{tab:mse-vs-turbo3}
\begin{tabular}{lrrrl}
  \toprule
  Model & \textsc{Turbo} K-MSE & turbo3 K-MSE & turbo3 / \textsc{Turbo} & Interpretation \\
  \midrule
  Qwen2.5-0.5B                  & $1.11\!\cdot\!10^{0}$  & $2.07\!\cdot\!10^{2}$  & $\mathbf{186\times}$   & we much better \\
  Qwen3-0.6B                    & $1.46\!\cdot\!10^{0}$  & $8.99\!\cdot\!10^{1}$  & $\mathbf{62\times}$    & we much better \\
  DeepSeek-R1-Distill           & $5.93\!\cdot\!10^{-1}$ & $1.44\!\cdot\!10^{3}$  & $\mathbf{2428\times}$  & we much better \\
  Gemma-4-E2B                   & $1.58\!\cdot\!10^{-3}$ & $2.96\!\cdot\!10^{-3}$ & $1.87\times$           & we better \\
  GLM-Edge-1.5B                 & $8.49\!\cdot\!10^{-1}$ & $2.81\!\cdot\!10^{-1}$ & $0.33\times$           & turbo3 better \\
  GLM-Edge-4B                   & $7.70\!\cdot\!10^{-1}$ & $1.48\!\cdot\!10^{-1}$ & $0.19\times$           & turbo3 better \\
  \bottomrule
\end{tabular}
\end{table}

On the Qwen/DeepSeek family (RoPE with halfsplit pairing), \textsc{KakeyaTurbo} has K-MSE $62\times$ to $\mathbf{2428\times}$ lower than turbo3. This gap corresponds exactly to the asymmetric-rescue configurations (q8\_0-K + turbo-V) that the TurboQuant documentation recommends for these architectures. On GLM-Edge (adjacent-RoPE + QK-norm), turbo3 has lower K-MSE because the K representation is already norm-stabilised by QK-norm, and the Gaussianised scalar quantization handles it cleanly. Gemma-4 is in between: \textsc{KakeyaTurbo} K-MSE is $1.87\times$ lower, reflecting the low-rank structure of Gemma's high-head-dim K.

\paragraph{Per-layer K-MSE distribution.}
Beyond the aggregate mean, the per-layer MSE distribution is informative because attention-critical layers may sit at distribution tails. Table~\ref{tab:mse-perlayer} gives the Qwen3-0.6B K-stream quantiles under v1.3 tier-1.

\begin{table}[H]
\centering
\small
\caption{Per-layer K-MSE distribution on Qwen3-0.6B under v1.3 tier-1, over 27 full-attention layers (L0 excluded).}
\label{tab:mse-perlayer}
\begin{tabular}{lr}
  \toprule
  Statistic & K-MSE \\
  \midrule
  min (layer 13) & $8.08\!\cdot\!10^{-1}$ \\
  25th percentile & $9.70\!\cdot\!10^{-1}$ \\
  median          & $1.14\!\cdot\!10^{0}$  \\
  75th percentile & $1.89\!\cdot\!10^{0}$  \\
  max (layer 4)   & $3.15\!\cdot\!10^{0}$  \\
  mean            & $1.46\!\cdot\!10^{0}$  \\
  \bottomrule
\end{tabular}
\end{table}

The max-to-min spread is $3.9\times$, smaller than the $200\times$ cross-layer spread that turbo3 exhibits on the same tensor set. This is a direct consequence of v1.3's per-block PCA: each layer fits its own low-rank structure, so worst-case layers do not diverge as dramatically as under a fixed rotation.

\subsection{Quality vs compression Pareto summary}

Ranking the codec configurations on the $(\text{ratio}, \text{quality})$ plane:
\begin{itemize}[leftmargin=*]
  \item Exact PCA with $b=3$ (v1.2 shipped default): low ratio, high quality.
  \item TurboQuant turbo3: mid ratio, low K quality on RoPE-dominated models (8$\times$--400$\times$ worse K MSE than v1.2 on the Qwen family per our prior measurements).
  \item \textsc{KakeyaTurbo} tier-1 ($b=2$, $r=D/2$): high ratio, MARGINAL quality on six models, ACCEPT on Gemma-4.
  \item Tier-1.5 (inverse-RoPE added): moves the Qwen / DeepSeek family further right by $+8\%$ to $+9\%$ ratio at no quality cost.
  \item SmolLM2 is the only model where no configuration enters the ACCEPT $\cap$ beats-turbo3 region (see \S\ref{sec:smollm2}).
\end{itemize}

\subsection{Flagship projections at 2026-Q2 scale}

None of the April-2026 flagship open-source models (Qwen3-235B-A22B, DeepSeek-V3.1, Kimi-K2-Instruct, GLM-4.6, MiniMax-M2) is loadable on our 15\,GB-RAM VM; their bf16 weights alone are $458$--$2000$\,GB. We therefore report byte-exact analytical extrapolation from architecturally-identical small-sibling proxies that \emph{are} measured (Table~\ref{tab:flagship}).

\begin{table}[H]
\centering
\small
\caption{Byte-exact extrapolation to 2026-Q2 flagship open-source models. Per-vector bytes come from the measured small-sibling proxy; $L$ and $n_{kv} \cdot \text{hd}$ come from the flagship config.}
\label{tab:flagship}
\begin{tabular}{lrrrrrr}
  \toprule
  Flagship & Params & $L$ & hd & tier-1 & tier-1.5 & vs turbo3 \\
  \midrule
  Qwen3-235B-A22B           & 235\,B & 94 & 128 & $6.53\times$ & $\mathbf{7.13\times}$ & $+39\%$ \\
  DeepSeek-V3.1 (decomp)    & 671\,B & 61 & 192 & $5.92\times$ & $6.76\times$ & $+30\%$ \\
  Kimi-K2-Instruct (decomp) & $1{,}000$\,B & 61 & 192 & $5.92\times$ & $6.76\times$ & $+30\%$ \\
  GLM-4.6                   & 355\,B & 92 & 128 & $5.84\times$ & ---              & $+14\%$ \\
  MiniMax-M2                & 229\,B & 62 & 128 & $5.92\times$ & $6.76\times$ & $+32\%$ \\
  \bottomrule
\end{tabular}
\end{table}

DeepSeek-V3.1 and Kimi-K2 use Multi-Latent Attention (MLA); the reported ratios are on the decompressed K/V (what the attention computation sees). The additional gain from applying \textsc{KakeyaTurbo} to the stored MLA latent is an open question.

\section{Related work}
\label{sec:related}

\paragraph{KV cache compression.} H\textsubscript{2}O~\cite{h2o} and Scissorhands~\cite{scissorhands} propose token eviction based on attention-score heavy hitters. StreamingLLM~\cite{streaming-llm} observes the attention-sink phenomenon and preserves the first few tokens exact. KIVI~\cite{kivi} is a per-channel quantization scheme for K (asymmetric quantization per channel) and per-token for V. TurboQuant~\cite{turboquant} is a data-oblivious WHT + Lloyd--Max scalar quantizer. We compare against turbo3 head-to-head on the same captured tensors in \S\ref{sec:benchmark}.

\paragraph{Low-rank KV methods.} Kakeya-KV~\cite{kakeya-kv} is our prior work and the direct ancestor of the skeleton construction here. DMC~\cite{dmc} and LESS~\cite{less} learn low-rank projections during training. Our approach is post-hoc and training-free.

\paragraph{Randomized numerical linear algebra.} The Halko--Martinsson--Tropp randomized SVD~\cite{halko-martinsson-tropp} is a textbook algorithm; we use it as a drop-in for exact PCA and observe that the implicit rank cap is the structural win. Frequent Directions~\cite{frequent-directions} is a related streaming PCA sketch we defer to future work.

\paragraph{Kakeya geometry and harmonic analysis.} The Kakeya problem originates in~\cite{kakeya-problem}. Its deep role in harmonic analysis began with Fefferman~\cite{fefferman-ball} (ball multiplier), continued through Bourgain's arithmetic-combinatorial bounds~\cite{bourgain-kakeya}, Wolff's geometric hairbrush argument~\cite{wolff-kakeya}, and Katz--{\L}aba--Tao's $\R^3$ improvements~\cite{katz-laba-tao}. The finite-field analogue was settled by Dvir's polynomial method~\cite{dvir-finite-field}; Guth's multilinear-Kakeya endpoint~\cite{guth-endpoint} introduced the tools eventually used by Wang and Zahl~\cite{wang-zahl} to resolve the Euclidean three-dimensional case. Our use of the Kakeya estimates is intuitional rather than a formal proof transfer: the scaling of the minimum tube-union measure (the common quantity across this literature) is the geometric sibling of the rate--distortion budget for a direction-covering code, and it motivates why the weighted PCA skeleton is the ``right'' dominant direction set to extract. Proposition~\ref{prop:kakeya-rd} states the analogy formally.

\section{Limitations and future work}
\label{sec:limitations}

\paragraph{Quality measurement.} We report reconstruction MSE, not perplexity or downstream-task accuracy. A full PPL sweep across LongBench/RULER is deferred to future work; early checks on Gemma-4-E2B at 2k context show greedy-decode token-level agreement with the bf16 baseline for the first 12 tokens. End-to-end quality at 128k context on flagship models remains unverified.

\paragraph{Flagship measurements.} All flagship numbers are analytical extrapolation from measured proxies. Validation on $\geq 500$\,GB-RAM hardware is pending.

\paragraph{MLA models.} DeepSeek-V3.1 and Kimi-K2 use MLA; our reported numbers compress the decompressed K/V, not the stored latent. Applying the codec to the MLA latent is a v1.4 open item.

\paragraph{GLM-family inverse-RoPE.} GLM-Edge uses adjacent-pairs RoPE + QK-norm; our halfsplit inverse is incompatible. A GLM-correct inverse rotation is tracked as v1.3.1.

\paragraph{Attention-aware extensions.} The \texttt{weights: \&[f32]} parameter in our API is a reserved socket for attention-score-driven importance, but is not yet wired. The v1.4 roadmap (public in the repository) proposes closing this loop by piping H\textsubscript{2}O-style prefill attention scores into the weighted PCA fit.

\paragraph{SmolLM2-class models.} MHA with head dim $64$ and flat V spectra defeats rank truncation; the Pareto gap observed on SmolLM2 will likely appear on any future model with similar structure. We ship SmolLM2 as a tier-2 capability that requires an explicit quality-versus-bytes trade-off at the deployment policy layer.

\section{Conclusion}
\label{sec:conclusion}

We presented \textsc{KakeyaTurbo}, a training-free KV cache compression codec that unifies TurboQuant-style scalar quantization and Kakeya-style low-rank skeleton construction within a single rate--distortion objective. Three design decisions --- (i) randomized-SVD skeleton with rank cap $r = D/2$, (ii) inverse-RoPE preprocessing on the key stream for canonical-RoPE models, and (iii) 2-bit Lloyd--Max residual quantization after WHT --- compose to a codec that beats TurboQuant turbo3 on 6 of 7 open-source models at ACCEPT-level quality, with measured gains of $+9.1\%$ to $+29.2\%$ at ctx = 4\,096 and analytical projections of $+14\%$ to $+39\%$ on 2026 flagship open-source models at 128k context. The project ships with 153 automated tests, real-data benchmarks on seven models, and a public repository with all raw measurements and runnable reproduction scripts. The honest position is: the codec is ready for deployment on Qwen, DeepSeek, Gemma, GLM-Edge, and MiniMax families; SmolLM2-class MHA-small-head models require a tier-2 quality-vs-bytes trade-off; end-to-end perplexity validation is future work; and the attention-aware extension (v1.4) has a concrete roadmap documented in the same repository.

\paragraph{Reproducibility.} All code, unit tests, integration tests, property-based tests, real-measurement JSON, and ablation decision documents are released under Apache-2.0. Repository:
\begin{center}
\url{https://github.com/FluffyAIcode/LLM-KV--Cache-compress}
\end{center}
The v1.3 release state lives on branch \texttt{cursor/v1-3-rsvd-rope-aware-12f5} (Pull Request \#11); the v1.4 roadmap is in \texttt{reports/v1\_3\_rsvd\_rope/V1\_4\_V1\_5\_ROADMAP.md}.

\begin{thebibliography}{99}

\bibitem{h2o}
Zhang, Z., Sheng, Y., Zhou, T., Chen, T., Zheng, L., Cai, R., Song, Z., Tian, Y., R\'e, C., Barrett, C., Wang, Z., Chen, B.
\newblock {H$_2$O: Heavy-hitter oracle for efficient generative inference of large language models.}
\newblock \emph{NeurIPS}, 2023.

\bibitem{scissorhands}
Liu, Z., Desai, A., Liao, F., Wang, W., Xie, V., Xu, Z., Kyrillidis, A., Shrivastava, A.
\newblock {Scissorhands: Exploiting the persistence of importance hypothesis for LLM KV cache compression at test time.}
\newblock \emph{NeurIPS}, 2023.

\bibitem{deepseekv3}
DeepSeek-AI.
\newblock {DeepSeek-V3 Technical Report.}
\newblock \emph{arXiv preprint arXiv:2412.19437}, 2024.

\bibitem{turboquant}
TheTom.
\newblock {TurboQuant+: Gaussian-bounded scalar quantization for KV cache.}
\newblock \emph{GitHub repository, \url{https://github.com/TheTom/turboquant_plus}}, 2024.

\bibitem{kakeya-kv}
KakeyaKV Contributors.
\newblock {Kakeya KV cache compression.}
\newblock \emph{\url{https://github.com/FluffyAIcode/LLM-KV--Cache-compress}}, 2026. (v1.0 and v1.2, this repository.)

\bibitem{kakeya-problem}
Kakeya, S.
\newblock {Some problems on maxima and minima regarding ovals.}
\newblock \emph{Tohoku Science Reports}, 6:71--88, 1917.

\bibitem{fefferman-ball}
Fefferman, C.
\newblock {The multiplier problem for the ball.}
\newblock \emph{Annals of Mathematics}, 94(2):330--336, 1971.

\bibitem{bourgain-kakeya}
Bourgain, J.
\newblock {Besicovitch type maximal operators and applications to Fourier analysis.}
\newblock \emph{Geometric and Functional Analysis}, 1(2):147--187, 1991.

\bibitem{wolff-kakeya}
Wolff, T.
\newblock {An improved bound for Kakeya type maximal functions.}
\newblock \emph{Revista Matem\'atica Iberoamericana}, 11(3):651--674, 1995.

\bibitem{katz-laba-tao}
Katz, N.~H., {\L}aba, I., Tao, T.
\newblock {An improved bound on the Minkowski dimension of Besicovitch sets in $\R^3$.}
\newblock \emph{Annals of Mathematics}, 152(2):383--446, 2000.

\bibitem{dvir-finite-field}
Dvir, Z.
\newblock {On the size of Kakeya sets in finite fields.}
\newblock \emph{Journal of the American Mathematical Society}, 22(4):1093--1097, 2009.

\bibitem{guth-endpoint}
Guth, L.
\newblock {The endpoint case of the Bennett-Carbery-Tao multilinear Kakeya conjecture.}
\newblock \emph{Acta Mathematica}, 205(2):263--286, 2010.

\bibitem{wang-zahl}
Wang, H., Zahl, J.
\newblock {Volume estimates for unions of convex sets, and the Kakeya set conjecture in three dimensions.}
\newblock \emph{arXiv preprint arXiv:2502.17655}, 2025.

\bibitem{halko-martinsson-tropp}
Halko, N., Martinsson, P.-G., Tropp, J.~A.
\newblock {Finding structure with randomness: Probabilistic algorithms for constructing approximate matrix decompositions.}
\newblock \emph{SIAM Review}, 53(2):217--288, 2011.

\bibitem{rope}
Su, J., Lu, Y., Pan, S., Murtadha, A., Wen, B., Liu, Y.
\newblock {RoFormer: Enhanced transformer with rotary position embedding.}
\newblock \emph{Neurocomputing}, 568:127063, 2024.

\bibitem{streaming-llm}
Xiao, G., Tian, Y., Chen, B., Han, S., Lewis, M.
\newblock {Efficient streaming language models with attention sinks.}
\newblock \emph{ICLR}, 2024.

\bibitem{kivi}
Liu, Z., Yuan, J., Jin, H., Zhong, S., Xu, Z., Braverman, V., Chen, B., Hu, X.
\newblock {KIVI: A tuning-free asymmetric 2bit quantization for KV cache.}
\newblock \emph{ICML}, 2024.

\bibitem{dmc}
Nawrot, P., Lancucki, A., Chochowski, M., Tarjan, D., Ponti, E.~M.
\newblock {Dynamic memory compression: Retrofitting LLMs for accelerated inference.}
\newblock \emph{arXiv preprint arXiv:2403.09636}, 2024.

\bibitem{less}
Dong, H., Yang, X., Zhang, Z., Wang, Z., Chi, Y., Chen, B.
\newblock {Get more with LESS: Synthesizing recurrence with KV cache compression for efficient LLM inference.}
\newblock \emph{ICML}, 2024.

\bibitem{frequent-directions}
Liberty, E.
\newblock {Simple and deterministic matrix sketching.}
\newblock \emph{KDD}, 2013.

\end{thebibliography}

\appendix

\section{Complete benchmark protocol}
\label{app:protocol}

All benchmarks in this paper are reproducible via the following one-line invocation after cloning the repository and downloading a model:
\begin{verbatim}
python3 benchmarks/kakeyaturbo_v1_2_real_bench.py \
    --model-path models/<model-dir> \
    --model-name <model-name> \
    --context-tokens 4096 \
    --bit-width 2 \
    --pca-method randomized \
    --out-dir reports/v1_3_rsvd_rope/bench/<model>_rsvd_half/ctx_4096
\end{verbatim}
The inverse-RoPE preprocessor adds:
\begin{verbatim}
python3 benchmarks/rope_aware_k_poc.py \
    --model-dir <model-dir> --model-name <model-name> \
    --ctx 4096 --rope-pairing halfsplit \
    --out-dir reports/v1_3_rsvd_rope/rope_poc/<model-name>
\end{verbatim}

Each run produces a JSON manifest per layer plus an aggregate \verb|summary.json|. All manifests are committed to the repository under \verb|reports/v1_3_rsvd_rope/|.

\section{Codec parameter defaults}
\label{app:params}

The reference v1.3 tier-1 configuration is:
\begin{center}
\begin{tabular}{ll}
  \toprule
  Parameter & Default value \\
  \midrule
  \texttt{bit\_width} & 2 \\
  \texttt{variance\_ratio} & 0.95 \\
  \texttt{k} (K-means centres) & 16 \\
  \texttt{block\_size} & 512 \\
  \texttt{kmeans\_max\_iter} & 32 \\
  \texttt{share\_basis} (V stream) & true \\
  \texttt{share\_basis} (K stream) & false \\
  \bottomrule
\end{tabular}
\end{center}

The randomized-SVD parameters are:
\begin{center}
\begin{tabular}{ll}
  \toprule
  Parameter & Default value \\
  \midrule
  \texttt{target\_rank} & $D/2$ \\
  \texttt{oversample} ($p$) & 8 \\
  \texttt{power\_iters} ($q$) & 2 \\
  \texttt{seed\_offset} & \texttt{0x9E37\_79B9\_7F4A\_7C15} \\
  \bottomrule
\end{tabular}
\end{center}

Tier-1.5 additionally applies the inverse-RoPE preprocessor to the K stream before the codec sees it, using the model's \texttt{rope\_theta} and halfsplit pairing convention.

\section{Qualitative ablation trail}
\label{app:ablation}

The v1.3 configuration is the endpoint of seven ablations run over the project's lifetime. For completeness, the full decision chain is:
\begin{enumerate}[leftmargin=*]
  \item \textbf{v1.0 $\to$ v1.2 A}: bf16 storage of skeleton (halves skeleton bytes, no quality loss). ACCEPT, shipped.
  \item \textbf{v1.2 B'}: V-stream layer-pooled PCA basis. ACCEPT (V only), shipped. Rejected on K due to family-dependent divergence.
  \item \textbf{K-side $d_{\text{eff}}$ truncation} + outlier compensation: every $0.05$ drop in \texttt{variance\_ratio} doubles K MSE; outliers recover $\sim 25\%$. \textbf{REJECT}.
  \item \textbf{Block size doubling} (512 $\to$ 1024): MSE inflation $1.14\times$ on average with $16\%$ byte savings. \textbf{MARGINAL}, not shipped as default.
  \item \textbf{Cross-layer basis sharing} on K: $2\times$--$15\times$ MSE inflation on Qwen/DeepSeek family, ACCEPT on Gemma/GLM-Edge. \textbf{REJECT as default, accept as per-family knob.}
  \item \textbf{Bit width sweep} ($b=3 \to b=2$): $+4\%$ to $+23\%$ ratio, MSE inflation $1.07\times$--$1.13\times$. \textbf{ACCEPT, new default.}
  \item \textbf{Randomized SVD with rank cap $r = D/2$}: combined with $b=2$, beats turbo3 on 6/7 at ACCEPT quality. \textbf{ACCEPT, new default.}
  \item \textbf{Inverse-RoPE K preprocessor}: $14\%$--$51\%$ K MSE drop on halfsplit-RoPE family, no-op or harmful elsewhere. \textbf{ACCEPT as per-family knob.}
\end{enumerate}
The full narrative is in the repository's \verb|reports/v1_3_rsvd_rope/DECISION.md|.

\end{document}
